\tablehead{tab:S5}{Physical Xiaomi Redmi Note~14 Pro+~5G mobile latency by
routing path during one session using a clean release build. All timing values
are reported in milliseconds.}
\begin{center}
\scriptsize
\setlength{\tabcolsep}{3pt}
\begin{tabular}{l c r r r r r r c c}
\toprule
      & Warm/ &      &        &     & Std- &     &     & Det.\ exec./ & Succ./ \\
Path  & meas. & Mean & Median & p95 & dev  & Min & Max & skipped      & fail.  \\
\midrule
Accepted Corn       & 10/100 & 147.73 & 148.58 & 157.59 & 10.26 & 137.33 & 207.18 & 100 / 0 & 100 / 0 \\
Accepted Pepper     & 10/100 & 151.81 & 151.29 & 160.91 & 8.61  & 139.42 & 192.18 & 100 / 0 & 100 / 0 \\
Accepted Tomato     & 10/100 & 150.25 & 150.70 & 158.67 & 5.80  & 139.05 & 163.42 & 100 / 0 & 100 / 0 \\
OOD rejection       & 10/100 & 53.64  & 53.12  & 57.58  & 2.07  & 52.33  & 63.94  & 0 / 100 & 100 / 0 \\
Crop mismatch       & 10/100 & 48.93  & 48.55  & 56.99  & 4.40  & 44.74  & 66.20  & 0 / 100 & 100 / 0 \\
\bottomrule
\end{tabular}
\end{center}

{\footnotesize
The five routing outcomes are defined in
the main-text protocol section.

\textit{Isolated-stage figures} (from a separate measured loop using the same
stage-timed functions, so they need not sum exactly to the end-to-end means;
this is run-to-run jitter, not an inconsistency): classifier inference mean
43.23~ms; image decode mean 0.69~ms; orientation-correction mean 0.079~ms;
preprocess mean 0.70~ms. Detector inference mean by input resolution:
$1920\times1080$, 93.75~ms; $1280\times960$, 92.79~ms; $640\times480$,
94.13~ms. Output-buffer transfer accounted for 0.54--0.55~ms in each case and was included
within detector inference. The Android processor-utilisation estimate was
0.98 and was derived from time spent by one processor thread; it is therefore
not directly comparable with the process-level estimate reported for iOS in
Table~\ref{tab:S4}.

\textit{Device and build.} Testing used a clean release build on a physical
Xiaomi Redmi Note~14 Pro+~5G running Android~16 with HyperOS. The device had a
QTI SM7635 processor, a 64-bit ARM architecture and approximately 11.09~GiB
of memory. The battery was at 100\% and charging throughout, while the device
reported its lowest thermal state. Charging is treated as part of the test
condition, and its effect is not estimated. The same classifier and detector
versions were used on both mobile platforms. The installed Android package
was approximately 381~MB, although this package-size estimate is not directly
comparable with the iOS bundle estimate. Inference used the CPU-based XNNPACK
backend; no mobile graphics, neural-network or vendor NPU backend was present.
The reported measurements must therefore not be described as NPU-accelerated.

\textit{Notes.} (1)~Memory samples from this run are not used for any claim
(see Table~\ref{tab:S7}) because the before and after readings were taken at
the same point rather than on opposite sides of the measured loop. This
sampling defect has since been corrected. (2)~Sub-millisecond
model-initialisation readings are excluded for the same reason as in
Table~\ref{tab:S4}. (3)~This is a single session
($n = 1$); run-to-run and thermal-drift variability on this device, as
observed for the iPhone in Table~\ref{tab:S4}, has not yet been
characterised.
\par}
\tabletail
